% GENERATED FILE. Edit canonical lessons, then run npm run generate:books.
% canonical-source-hash: fnv1a64:7d5e3dde6a1a78f1
% canonical-lessons: LA-C415-castigo, LA-C415-relego, LA-C415-proscribo, LA-C415-flagellum, LA-C415-custodia

\chapter{Punishment}
\label{ch:la-punishment}
\hlchaptermodality{415}

\begin{chapteropening}
\textbf{By the end of this chapter:} \emph{I can say that someone was punished, banished or outlawed, and talk about the whip and custody.}

Complete the last lesson of chapter 415: I can say that someone was punished, banished or outlawed, and talk about the whip and custody.
\end{chapteropening}

\section[castīgō, castīgāre]{castīgō, castīgāre — to punish, to correct; to rebuke}

\label{lesson:LA-C415-castigo}



\begin{quote}
Before the new one: say the Latin for a crime against heaven, then the Latin for unfair.
\end{quote}

\subsection*{You'll want to know: castīgō, castīgāre}
\textbf{castīgō, castīgāre} — \textquotedblleft{}to punish, to correct; to rebuke\textquotedblright{}.

\begin{quote}
Pater fīlium verbīs castīgāvit.
\end{quote}

\begin{quote}
\emph{The father rebuked his son.}
\end{quote}

\emph{castīgō, castīgāre, castīgāvī, castīgātum}, built on \emph{castus}, pure: to set right. English \emph{castigate} and \emph{chastise}.

\begin{cousinweb}
\emph{castīgō} is built on \emph{castus} \textquotedblleft{}pure\textquotedblright{} (EtymDB; Lewis \& Short): to set a person right. English \emph{castigate} and \emph{chastise} come from it.
\end{cousinweb}

\begin{grammarlens}[title={What you've built}]
One more word for telling, explaining and describing at B1.
\end{grammarlens}

\subsection*{Your turn}
\begin{itemize}
\raggedright
  \item \emph{Say it:} \emph{castīgō, castīgāre}
  \item \emph{Say it:} \emph{castīgō, castīgāre}, once more
  \item \emph{Say it:} a sentence of your own with \emph{castīgō, castīgāre}
  \item \emph{Recall:} say the Latin for a crime against heaven, then the Latin for unfair, then say \emph{castīgō, castīgāre} again
\end{itemize}

\begin{tcolorbox}[breakable,colback=teal!4,colframe=teal!35!black,title={Before you move on}]
What does castīgō, castīgāre mean? (\textquotedblleft{}To punish, to correct; to rebuke\textquotedblright{}.) Say it once more.
\end{tcolorbox}
\section[relēgō, relēgāre]{relēgō, relēgāre — to send away, to banish}

\label{lesson:LA-C415-relego}



\begin{quote}
Before the new one: say the Latin for unfair, then the Latin for to punish.
\end{quote}

\subsection*{You'll want to know: relēgō, relēgāre}
\textbf{relēgō, relēgāre} — \textquotedblleft{}to send away, to banish\textquotedblright{}.

\begin{quote}
Prīnceps eum procul ab urbe relēgāvit.
\end{quote}

\begin{quote}
\emph{The emperor banished him far from the city.}
\end{quote}

\emph{relēgō, relēgāre, relēgāvī, relēgātum}: a milder exile, which kept a man's property and citizenship. English \emph{relegate}.

\begin{cousinweb}
\emph{relēgō} is built from \emph{re-} \textquotedblleft{}away, back\textquotedblright{} and \emph{lēgō} \textquotedblleft{}to send with a charge, to depute\textquotedblright{} (EtymDB; Lewis \& Short). English \emph{relegate} comes from it.
\end{cousinweb}

\begin{grammarlens}[title={What you've built}]
One more word for telling, explaining and describing at B1.
\end{grammarlens}

\subsection*{Your turn}
\begin{itemize}
\raggedright
  \item \emph{Say it:} \emph{relēgō, relēgāre}
  \item \emph{Say it:} \emph{relēgō, relēgāre}, once more
  \item \emph{Say it:} a sentence of your own with \emph{relēgō, relēgāre}
  \item \emph{Recall:} say the Latin for unfair, then the Latin for to punish, then say \emph{relēgō, relēgāre} again
\end{itemize}

\begin{tcolorbox}[breakable,colback=teal!4,colframe=teal!35!black,title={Before you move on}]
What does relēgō, relēgāre mean? (\textquotedblleft{}To send away, to banish\textquotedblright{}.) Say it once more.
\end{tcolorbox}
\section[prōscrībō, prōscrībere]{prōscrībō, prōscrībere — to post for sale; to outlaw, to proscribe}

\label{lesson:LA-C415-proscribo}



\begin{quote}
Before the new one: say the Latin for to punish, then the Latin for to send away.
\end{quote}

\subsection*{You'll want to know: prōscrībō, prōscrībere}
\textbf{prōscrībō, prōscrībere} — \textquotedblleft{}to post for sale; to outlaw, to proscribe\textquotedblright{}.

\begin{quote}
Multī cīvēs prōscrīptī sunt.
\end{quote}

\begin{quote}
\emph{Many citizens were outlawed.}
\end{quote}

\emph{prōscrībō, prōscrībere, prōscrīpsī, prōscrīptum}: to post a notice, first of a house for sale, then of a man whose life and goods were forfeit. English \emph{proscribe}.

\begin{cousinweb}
\emph{prōscrībō} is built from \emph{prō} \textquotedblleft{}forth, in public\textquotedblright{} and \emph{scrībō} \textquotedblleft{}to write\textquotedblright{} (EtymDB; Lewis \& Short): to post a notice. English \emph{proscribe} comes from it.
\end{cousinweb}

\begin{grammarlens}[title={What you've built}]
One more word for telling, explaining and describing at B1.
\end{grammarlens}

\subsection*{Your turn}
\begin{itemize}
\raggedright
  \item \emph{Say it:} \emph{prōscrībō, prōscrībere}
  \item \emph{Say it:} \emph{prōscrībō, prōscrībere}, once more
  \item \emph{Say it:} a sentence of your own with \emph{prōscrībō, prōscrībere}
  \item \emph{Recall:} say the Latin for to punish, then the Latin for to send away, then say \emph{prōscrībō, prōscrībere} again
\end{itemize}

\begin{tcolorbox}[breakable,colback=teal!4,colframe=teal!35!black,title={Before you move on}]
What does prōscrībō, prōscrībere mean? (\textquotedblleft{}To post for sale; to outlaw, to proscribe\textquotedblright{}.) Say it once more.
\end{tcolorbox}
\section[flagellum, flagellī]{flagellum, flagellī — a whip, a scourge}

\label{lesson:LA-C415-flagellum}



\begin{quote}
Before the new one: say the Latin for to send away, then the Latin for to post for sale.
\end{quote}

\subsection*{You'll want to know: flagellum, flagellī}
\textbf{flagellum, flagellī} — \textquotedblleft{}a whip, a scourge\textquotedblright{}.

\begin{quote}
Aurīga flagellum in manū tenet.
\end{quote}

\begin{quote}
\emph{The charioteer holds a whip in his hand.}
\end{quote}

\emph{flagellum, flagellī} n., a small form of \emph{flagrum}, a whip. English \emph{flagellate}.

\begin{cousinweb}
\emph{flagellum} is a small form of \emph{flagrum} \textquotedblleft{}a whip, a scourge\textquotedblright{} (EtymDB; Lewis \& Short). English \emph{flagellate} comes from it.
\end{cousinweb}

\begin{grammarlens}[title={What you've built}]
One more word for telling, explaining and describing at B1.
\end{grammarlens}

\subsection*{Your turn}
\begin{itemize}
\raggedright
  \item \emph{Say it:} \emph{flagellum, flagellī}
  \item \emph{Say it:} \emph{flagellum, flagellī}, once more
  \item \emph{Say it:} a sentence of your own with \emph{flagellum, flagellī}
  \item \emph{Recall:} say the Latin for to send away, then the Latin for to post for sale, then say \emph{flagellum, flagellī} again
\end{itemize}

\begin{tcolorbox}[breakable,colback=teal!4,colframe=teal!35!black,title={Before you move on}]
What does flagellum, flagellī mean? (\textquotedblleft{}A whip, a scourge\textquotedblright{}.) Say it once more.
\end{tcolorbox}
\section[custōdia, custōdiae]{custōdia, custōdiae — custody, guard; prison}

\label{lesson:LA-C415-custodia}



\begin{quote}
Before the new one: say the Latin for to post for sale, then the Latin for a whip.
\end{quote}

\subsection*{You'll want to know: custōdia, custōdiae}
\textbf{custōdia, custōdiae} — \textquotedblleft{}custody, guard; prison\textquotedblright{}.

\begin{quote}
Reus in custōdiā erat.
\end{quote}

\begin{quote}
\emph{The accused was in custody.}
\end{quote}

\emph{custōdia, custōdiae} f., built on \emph{custōs}, a guard. \emph{In custōdiam dare}, to put under guard. English \emph{custody}.

\begin{cousinweb}
\emph{custōdia} is built on \emph{custōs, custōdis} \textquotedblleft{}a guard\textquotedblright{} (EtymDB; Lewis \& Short). English \emph{custody} comes from it.
\end{cousinweb}

\begin{grammarlens}[title={What you've built}]
One more word for telling, explaining and describing at B1.
\end{grammarlens}

\subsection*{Your turn}
\begin{itemize}
\raggedright
  \item \emph{Say it:} \emph{custōdia, custōdiae}
  \item \emph{Say it:} \emph{custōdia, custōdiae}, once more
  \item \emph{Say it:} a sentence of your own with \emph{custōdia, custōdiae}
  \item \emph{Recall:} say the Latin for to post for sale, then the Latin for a whip, then say \emph{custōdia, custōdiae} again
\end{itemize}

\begin{tcolorbox}[breakable,colback=teal!4,colframe=teal!35!black,title={Before you move on}]
What does custōdia, custōdiae mean? (\textquotedblleft{}Custody, guard; prison\textquotedblright{}.) Say it once more.
\end{tcolorbox}
